\chapter*{Prakata}\pagestyle{preface}\thispagestyle{preface}
\setlength{\parskip}{.25ex}


Kumpulan ini melengkapi buku \nocite{Hefferon12}
\textit{Aljabar Linear}.\footnote{Lihat
\protect\url{hefferon.net/linearalgebra}
untuk PDF yang dapat diunduh secara bebas, materi pendamping, dan
sumber \protect\LaTeX{}.}
Kumpulan ini membantu mahasiswa memperkuat dan memperluas pemahaman mereka
tentang aljabar linear dengan menggunakan perangkat lunak matematika
\Sage{}.% \footnote{See 
% \url{http://www.sagemath.org} for the software and documentation.}

Salah satu tujuan utama setiap program sarjana Matematika adalah membawa
mahasiswa menuju pemahaman yang lebih tinggi tentang bidang tersebut.
Pada awal studi, mereka mengikuti mata kuliah Kalkulus yang mengurangi
kelengkapan dan keketatan demi latihan algoritmik, sedangkan mata kuliah
selanjutnya mencurahkan lebih banyak perhatian pada konsep dan pembuktian.
\textit{Aljabar Linear} sesuai dengan rencana perkembangan ini.
Buku tersebut berupaya membawa mahasiswa pada pemahaman yang lebih dalam,
sekaligus mengakui bahwa pada tahap ini cukup banyak perhitungan membantu
proses tersebut dengan memperkuat konsep-konsepnya.

Wajar jika buku menggunakan contoh dan soal pekerjaan rumah yang kecil dengan
bilangan yang mudah dikelola\Dash tugas mengalikan dua matriks tiga kali tiga
berentri bilangan bulat kecil dengan tangan akan membangun intuisi, sedangkan
meminta mahasiswa melakukan hal yang sama pada matriks dua puluh kali dua
puluh dengan bilangan sepuluh angka desimal hanya akan menyiksa mereka.
Namun, seorang pengajar dapat khawatir bahwa pendekatan ini melewatkan
kesempatan untuk memperdalam pemahaman mahasiswa melalui eksplorasi yang tidak
dibatasi kertas dan pensil, atau untuk menunjukkan bahwa bidang ini diterapkan
secara luas, sebab perhitungan dalam contoh penerapan realistis biasanya
sulit.

Perangkat lunak matematika dapat memperluas apa yang wajar dikerjakan hingga
mencakup sistem yang lebih besar, bilangan yang lebih rumit, dan perhitungan
yang lebih panjang.
Manual ini memperluas kemampuan mahasiswa ke arah tersebut.
Selain membantu memperkuat pemahaman konsep, mempelajari cara mengotomatiskan
pekerjaan dan menangani tugas yang lebih besar membuat pengalaman mahasiswa
lebih menyerupai pekerjaan seorang profesional dalam praktik.
Keuntungan lain adalah mahasiswa menjumpai gagasan baru, seperti ukuran
pertumbuhan waktu jalan.

Kalau begitu, mengapa tidak langsung mengajar dengan komputer?

Tujuan kita adalah mengembangkan pemahaman yang lebih tinggi tentang bidang
ini. Untuk itu, kita mempertahankan fokus pada ruang vektor dan pemetaan
linear. Uraian di sini memperlakukan komputasi sebagai alat untuk membangun
pemahaman tersebut, bukan sebagai objek utama.

Sebagian pengajar mendapati bahwa mahasiswa mereka paling terbantu dengan
tetap berpegang pada materi inti.
Pengajar lain memiliki mahasiswa yang akan memperoleh manfaat dari jenis
pekerjaan yang kita lakukan di sini.
Karena manual ini merupakan pendamping, pengajar yang berbeda dapat membuat
pilihan yang berbeda pula.


\section{Mengapa \Sage?}
Dalam
\textit{Open Source Mathematical Software}
\citep{JoynerStein07},\footnote{Teks lengkap tersedia di
\protect\url{www.ams.org/notices/200710/tx071001279p.pdf}.}
para penulis berpendapat bahwa jalan terbaik bagi Matematika adalah
memperhatikan lisensi perangkat lunak yang kita gunakan.

\begin{quotation}\small
Andaikan Jane adalah matematikawan terkenal yang mengumumkan bahwa ia telah
membuktikan suatu teorema. Kita mungkin akan memercayainya, tetapi ia tahu
bahwa ia harus memberikan pembuktian jika diminta. Sekarang andaikan Jane
menyatakan suatu teorema benar dengan sebagian dasar berupa hasil perangkat
lunak. Tanpa gagasan baru, hal terdekat yang secara masuk akal dapat kita
harapkan dari pembuktian yang ketat adalah kemampuan untuk memeriksa secara
terbuka dan menggunakan seluruh kode komputer yang mendasari hasil tersebut.
Jika programnya berpemilik, hal ini tidak mungkin. Kita sepenuhnya berhak
bersikap tidak percaya, bukan hanya karena ketidakpercayaan samar terhadap
komputer, tetapi karena pemrogram terbaik pun rutin melakukan kesalahan.

Jika seseorang membaca pembuktian teorema Jane dengan harapan memperluas
gagasannya atau menerapkannya dalam konteks baru, tidak memiliki akses ke
cara kerja internal perangkat lunak yang mendasari hasil Jane merupakan
suatu batasan.
\end{quotation}  
Para profesional memilih alat dengan menyeimbangkan banyak faktor, tetapi
argumen ini meyakinkan.
Di sini kita memakai \Sage{} karena kemampuannya yang tinggi dan karena
perangkat lunak ini Bebas, yang berarti bukan sekadar tersedia tanpa biaya,
melainkan juga bahwa kode sumbernya terbuka untuk diperiksa oleh siapa pun.
\footnote{Lihat \protect\url{www.gnu.org/philosophy/free-sw.html}
untuk latar belakang dan definisinya.}



\section{Manual ini}
Manual ini Bebas.
Lihat \url{hefferon.net/linearalgebra} untuk lisensinya.
Secara khusus, mahasiswa bebas mengunduhnya dan pengajar bebas
mendistribusikannya kembali, misalnya melalui situs intranet mata kuliah.

%% I am glad to get feedback, especially from instructors
%% who have class tested the material.
%% That same link gives my current contact information.

Sistem perangkat lunak yang kita gunakan untuk menjelajahi bidang ini,
\python{} dan \Sage{}, juga Bebas.
Pengguna mungkin sudah memilikinya pada sistem operasi masing-masing;
jika belum, keduanya tersedia dari
\href{http://www.python.org}{\url{www.python.org}}
dan
\href{https://www.sagemath.org}{\url{www.sagemath.org}}.
Keduanya merupakan sistem perangkat lunak yang kuat dan cakap, yang telah
digunakan dalam lingkungan produksi oleh ribuan profesional selama bertahun-tahun.

\textbf{Catatan tentang ketepatan contoh:}~kita sering menampilkan kode
masukan untuk \python{} dan \Sage{} beserta tanggapan komputer.
Tampilan tersebut tidak disalin-tempel.
Sebaliknya, interaksi itu direkam sebagai bagian dari proses pembuatan
dokumen ini. Karena itu, yang Anda lihat seharusnya tepat, kecuali versi
perangkat lunak Anda sangat berbeda dari versi yang saya gunakan.
Berikut versi \Sage{} saya.
\begin{sagecommandline}
sage: version()  
\end{sagecommandline}
Berikut versi \python{} saya.
\begin{pythonconsole}
import sys
print(sys.version)
\end{pythonconsole}




\section{Membaca manual ini}
Karena manual ini melengkapi buku \textit{Aljabar Linear}, kita tidak
mendefinisikan istilah atau membuktikan hasil di sini.
Sebaliknya, mahasiswa sebaiknya mempelajari materi ini setelah menyelesaikan
bab terkait dalam buku dan memakai bab tersebut sebagai rujukan.

Kaitan antara materi di sini dan buku adalah sebagai berikut:
\textit{Python dan Sage} tidak bergantung pada buku;
\textit{Metode Gauss} berkaitan dengan Bab Satu;
\textit{Ruang Vektor} untuk Bab Dua;
\textit{Matriks}, \textit{Pemetaan}, dan
\textit{Dekomposisi Nilai Singular} berkaitan dengan Bab Tiga;
\textit{Geometri Pemetaan Linear} paling sesuai dengan Bab Empat; dan
\textit{Nilai Eigen} sesuai dengan Bab Lima
(bagian itu menyebut Bentuk Jordan, tetapi hanya bergantung pada materi
sampai diagonalisasi).

%% An instructor may want to write Jupyter notebooks for the chapters here
%% since they do not have stand-alone exercises.




\section{Ucapan Terima Kasih}
Saya senang mendapat kesempatan untuk berterima kasih kepada tim pengembang
\python{} dan \Sage{}.
Secara khusus, karya ini tidak akan terwujud tanpa \citep{SageTeam19ref}.
Saya juga senang dapat menyebut \citep{Beezer11} sebagai sumber inspirasi.
Terakhir, saya berterima kasih kepada Saint Michael's College atas waktu yang
diberikan untuk menulis dokumen ini.





\vspace{.5in}
\noindent\begin{tabular}[t]{@{}l}
\textit{Kami menekankan latihan.} \\
\quotefrom{Suzuki}  \\[2ex]
\textit{Penyajian yang teratur tidak selalu buruk} \\
\quad\textit{tetapi mungkin tidak memadai jika berdiri sendiri.} \\
\quotefrom{Brandt}  
\end{tabular}

\vspace*{.25in}
\hbox{}
\hfill
\begin{tabular}[t]{l@{}}
Jim Hef{}feron \\
Matematika, Saint Michael's College \\
Colchester, Vermont, AS \\
2021-Oct-12
\end{tabular}
\vspace{1ex}
\endinput

TODO
